\section{The original lower-root projection: exact action, zero kernel, and a finite determinant formula with o(kq) error}\label{the-original-lower-root-projection-exact-action-zero-kernel-and-a-finite-determinant-formula-with-okq-error}

21 September 2026. This continuation treats the actual matrices \(S_N\) remaining in the arithmetic receiver. It preserves every original lower root, the physical center and phase, all four source degrees, and the attained quotient projection. The result is an explicit finite action on the original pole-canceling functions, with a proved logarithmic determinant error of order \(O(k)\) on that row space. The complementary pole directions contribute the stated \(o(kq)\) bound. No value or sign for the resulting determinant is asserted without evaluating its displayed matrix.

\subsection{1. The complete original projection and its polynomial kernel}\label{the-complete-original-projection-and-its-polynomial-kernel}

Use all of KF1--55, CGR1--29 and CGS1--22 with their original domains. In addition impose \(k\ge17\), which makes all radii below positive. At an original cutoff \(N\in\{q-1,q,2q-1,2q\}\), put \[
D=N-v,\quad M=D-q'=N-q+g,\qquad
I_N=\chi'\mathcal P'_{M}\subset\mathcal P'_D.
\tag{CGP1}
\] The four values of \(M\) are exactly \(g-1,g,q+g-1,q+g\). The original lower roots and their physical coordinates are \[
\beta'_{ab}=c'+(2a-k')\delta+i(2b-k')\gamma,\qquad
\omega'_{ab}=(\beta'_{ab}-c')/i
=(2b-k')\gamma-i(2a-k')\delta,
\tag{CGP2}
\] where \(k'=k-8\), \(c'=k/2-4\), and every \(0\le a,b\le k'\) occurs. These \(q'\) roots are pairwise distinct. The complete physical lower polynomial is \(i^{-q'}\chi'(c'+iy)=\prod_{a,b}(y-\omega'_{ab})\).

Let \(p_n\) be the monic original Gamma polynomials of CG8, and retain their exact squared norms \[
\gamma_n=M_\sigma n!(1/2)_n,\qquad M_\sigma=\sqrt{2\pi},
\quad \phi_{n,c'}(S')=p_n((S'-c')/i)/\sqrt{\gamma_n}.
\tag{CGP3}
\] Their orthonormal coefficient frame identifies \(\mathcal P'_D\) with \(\mathbb C^{D+1}\) isometrically. Define the full evaluation matrix and its exact positive Gram by \[
(E_D)_{\alpha n}=p_n(\omega'_\alpha)/\sqrt{\gamma_n},
\quad 0\le n\le D,\qquad G_D=E_DE_D^*.
\tag{CGP4}
\] This matrix has full row rank: \(D\ge q'-1\), and the usual degree-\((q'-1)\) Lagrange polynomials interpolate all the displayed distinct roots. Consequently \(G_D>0\). The kernel of evaluation is exactly \(I_N\). Indeed a polynomial in \(\mathcal P'_D\) vanishes at all these roots if and only if it is divisible by their entire polynomial \(\chi'\); the remaining degree is at most \(D-q'=M\). The exact orthogonal projection is therefore \[
\boxed{P_{I_N}=I_{D+1}-E_D^*G_D^{-1}E_D.}
\tag{CGP5}
\] Direct multiplication shows that the right side is self-adjoint and idempotent, has image \(\ker E_D\), and kills the orthogonal complement \(\operatorname{im}E_D^*\). This proves the formula with every original lower-root constraint present.

The entries of \(G_D\) have the exact Christoffel--Darboux expression \[
K_D(x,z)=\sum_{n=0}^D\frac{p_n(x)p_n(\bar z)}{\gamma_n}
=\frac{p_{D+1}(x)p_D(\bar z)-p_D(x)p_{D+1}(\bar z)}
{\gamma_D(x-\bar z)},\qquad
(G_D)_{\alpha\beta}=K_D(\omega'_\alpha,\omega'_\beta).
\tag{CGP6}
\] At \(x=\bar z\), use the polynomial limit \([p'_{D+1}(x)p_D(x)-p'_D(x)p_{D+1}(x)]/\gamma_D\). For a proof, the exact recurrence is \(p_{n+1}(x)=xp_n(x)-n(n-1/2)p_{n-1}(x)\), with \(\gamma_n=n(n-1/2)\gamma_{n-1}\). Multiply the recurrence at \(x\) and \(\bar z\) by the opposite polynomial and subtract. After division by \(\gamma_n\), the resulting differences telescope from \(n=0\) to \(D\), giving CGP6. No estimate of a Vandermonde determinant or replacement of the complete lower polynomial is used.

\subsection{2. Exact action on the original invariant functions, and its kernel}\label{exact-action-on-the-original-invariant-functions-and-its-kernel}

For an original invariant row \(z\) after the full low constraint, retain exactly \[
F_z(w)=\sum_{\alpha=1}^{q}z_\alpha e^{\beta_\alpha w},\qquad
G_z(w)=F_z(w)/E_A(w),\qquad
w(t)=-i\arctan t,
\] \[
H_z(t)=(1+t^2)^{-1/4}e^{-c'w(t)}G_z(w(t))
=\sum_{n\ge0}a_{z,n}t^n.
\tag{CGP7}
\] The function is analytic at zero since \(z\) kills all polynomials of degree below \(v=\operatorname{ord}_0E_A\). Its original functional row in the lower Gamma frame is \[
W_{z,N}=(\rho_na_{z,n}/\sqrt{M_\sigma})_{n=0}^D,
\qquad \rho_n=\sqrt{n!/(1/2)_n}.
\tag{CGP8}
\] This is exactly CGR19, obtained by applying the polynomial generating function to \(G_z\); in particular it retains \(c'\) and the phase \(-i\) in \(w(t)\). Equations CGP5 and CGP8 now give the explicit original action \[
\boxed{W_{z,N}P_{I_N}
=W_{z,N}-(W_{z,N}E_D^*)G_D^{-1}E_D.}
\tag{CGP9}
\] For two actual rows \(z,z'\), its Gram entry is exactly \[
W_{z,N}W_{z',N}^*
-(W_{z,N}E_D^*)G_D^{-1}(E_DW_{z',N}^*).
\tag{CGP10}
\] Thus CGP6--10 specify the action of the complete lower-root projection on the original functions by finite polynomial sums and a fixed \(q'\)-row Gram inverse.

This projected row map has zero kernel on the actual invariant row space at every original cutoff. Here is a coefficient proof that does not infer this from its dimension. If \(W_{z,N}P_{I_N}=0\), CGP5 implies \(W_{z,N}=bE_D\) for some full lower-root coefficient row \(b=(b_\alpha)\). Equality of these Gamma functional coefficients is equality on every polynomial of degree at most \(D\). Hence the first \(D+1\) Taylor coefficients at zero agree for \[
G_z(w)\quad\hbox{and}\quad B_b(w)=\sum_{\alpha=1}^{q'}
b_\alpha e^{\beta'_\alpha w}.
\] The difference has zero of order at least \(D+1\). Multiplication by \(E_A\), of order exactly \(v\), therefore gives \[
F_z(w)-E_A(w)B_b(w)=O(w^{N+1}).
\tag{CGP11}
\] Root addition \(\beta'_{ij}+b_{rs}=\beta_{i+r,j+s}\) identifies the second term with the original conductor coefficient vector \(M_Ab\) on the entire original \(q\)-root set. Since \(N\ge q-1\), its first \(q\) moments determine that vector uniquely: the Vandermonde evaluation matrix on these distinct roots is invertible by their Lagrange interpolation polynomials. Thus \(z=M_Ab\). But \(z\) is supported on the actual strip, and OCF7 gives \(N_kz=z\), \(N_kM_A=0\). Therefore \(z=0\). This proves the zero kernel by the exact original conductor map. It does not assert a uniform lower singular value from injectivity.

\subsection{3. Analytic control on the actual pole-canceling row space}\label{analytic-control-on-the-actual-pole-canceling-row-space}

Take precisely \(\mathcal W_k^0\) of CGR23. Its rows cancel every specified value and jet at every zero of \(g_A(t)=t^{-v}e^{-4w(t)}E_A(w(t))\) in \(|t|<1-k^{-1/2}\). Let its dimension be \(t_k=r-h_k\), with \(h_k\le2J_g\sqrt k\). Set \[
\eta=k^{-1/2},\quad s=1-2\eta,\quad r_c=1-3\eta,
\quad 0<r_c<s<1-\eta.
\tag{CGP12}
\] Use exactly CGR21 and CGR25's finite constants \(J_g\) and \(F_k\). Their Blaschke/Harnack proof gives, for every such original row, \[
\max_{|t|\le s}|H_z(t)|\le A_k\|z\|_2,
\qquad A_k=F_k\exp(2J_g/\eta).
\tag{CGP13}
\] Every canceled zero is interpreted by its analytic removable value; CGP13 does not evaluate an indeterminate ratio as zero. In particular \(H_z\) is holomorphic on a neighborhood of this closed disk. The proof uses the actual pole cancellations, not an assumed absence of conductor zeros.

Choose a fixed coefficient-orthonormal basis \(z_1,\ldots,z_{t_k}\) of \(\mathcal W_k^0\), so \(Z_0Z_0^*=I_{t_k}\). This is only a stated exact change of the row frame, shared at all four cutoffs; no source norm is altered. Let \(W_N^0\) have rows CGP8 in this basis. From CGP13, Cauchy's coefficient bound and \(\rho_n^2\le2n+1\), \[
\|W_N^0\|^2\le B_N^2,
\qquad
B_N^2=\max\left\{1,
\frac{(D+1)(2D+1)}{M_\sigma}A_k^2s^{-2D}\right\}.
\tag{CGP14}
\] This bound is uniform on combinations of the actual rows because their coefficient norm is exactly the norm of their combination coefficients. It is an operator-norm bound, without a row-count factor. CGR27 gives \(\log B_N^2=o(q)\).

The complete projected covariance lower bound CGS6 remains valid on this row space. Put \[
\lambda_N=\min\left\{1,
\frac1{M_D^2B_{\mathrm{val}}q
(C_A^{\mathrm{path}})^{20k}}\right\}>0,
\tag{CGP15}
\] where the product in the denominator means \(M_D^2\,B_{\mathrm{val}}\,q\,(C_A^{\mathrm{path}})^{20k}\), with all constants from CGS1--6. Then \[
\lambda_N I\preceq C_N^0:=W_N^0P_{I_N}(W_N^0)^*
\preceq\widetilde C_N^0:=W_N^0(W_N^0)^*,\qquad
\log(1/\lambda_N)=O_A(q\log(q+2)).
\tag{CGP16}
\] The minimum with one is a smaller proved lower constant, used only to make the finite estimates below immediate. The full original metric and the projection CGP5 remain unchanged.

\subsection{4. A finite Cauchy action with an explicit error}\label{a-finite-cauchy-action-with-an-explicit-error}

For an integer \(T>D\), write \(\omega_j=e^{2\pi ij/T}\) and define coefficients using actual values of the original function: \[
a_{z,n}^{[T]}=\frac{r_c^{-n}}{T}
\sum_{j=0}^{T-1}H_z(r_c\omega_j)\omega_j^{-n},
\qquad 0\le n\le D.
\tag{CGP17}
\] Let \(W_N^{[T]}\) be the matrix obtained from CGP8 using these coefficients and the same rows \(z_1,\ldots,z_{t_k}\). This is a finite formula; all its source degrees remain \(0,\ldots,D\). It is not a truncation of the lower root polynomial.

Expanding \(H_z\) on the radius-\(r_c\) circle and summing the finite geometric progression proves the exact alias identity \[
a_{z,n}^{[T]}-a_{z,n}
=\sum_{j\ge1}a_{z,n+jT}r_c^{jT}.
\tag{CGP18}
\] There are no negative alias indices because \(0\le n\le D<T\). The series is absolutely convergent by CGP13. Cauchy's coefficient bound therefore gives, with \(\tau=(r_c/s)^T\), \[
|a_{z,n}^{[T]}-a_{z,n}|
\le A_k\|z\|_2s^{-n}\frac{\tau}{1-\tau}.
\] Applying this to every row combination and using exactly the same sum as CGP14 gives \[
\boxed{\|W_N^{[T]}-W_N^0\|
\le B_N\frac{\tau}{1-\tau}.}
\tag{CGP19}
\] The error has no unestimated Gram inverse or condition number.

Choose the explicit integer \[
T_N=\max\left\{D+1,
\left\lceil\frac{\log(32B_N^2/\lambda_N)}
{\log(s/r_c)}\right\rceil\right\}.
\tag{CGP20}
\] Then \(\tau\le\lambda_N/(32B_N^2)\le1/32\), and the operator error \(e_N\) in CGP19 obeys \[
e_N\le\lambda_N/(31B_N).
\tag{CGP21}
\] Since \(B_N\ge1\) and \(\lambda_N\le1\), \(2B_Ne_N+e_N^2\le(2/31+1/31^2)\lambda_N
<\lambda_N/4\). Set, using the exact original projection, \[
C_N^{[T]}=W_N^{[T]}P_{I_N}(W_N^{[T]})^*,\qquad
\widetilde C_N^{[T]}=W_N^{[T]}(W_N^{[T]})^*.
\tag{CGP22}
\] For either covariance, expanding its two factors and using \(\|P_{I_N}\|=1\) proves that its error has norm at most \(2B_Ne_N+e_N^2<\lambda_N/4\). By CGP16 this is at most one quarter of the corresponding exact covariance in positive-form order. Consequently \[
\frac34 C_N^0\preceq C_N^{[T]}\preceq\frac54 C_N^0,
\qquad
\frac34\widetilde C_N^0\preceq\widetilde C_N^{[T]}
\preceq\frac54\widetilde C_N^0.
\tag{CGP23}
\] Both finite matrices are strictly positive. Define their exact finite determinant ratio \[
\mathcal A_N^{[T]}=
\log\frac{\det C_N^{[T]}}{\det\widetilde C_N^{[T]}},
\qquad
S_N^0=(\widetilde C_N^0)^{-1/2}C_N^0
(\widetilde C_N^0)^{-1/2}.
\] Taking logarithmic determinants in CGP23 gives the quantitative projection-action conclusion \[
\boxed{|\log\det S_N^0-\mathcal A_N^{[T]}|
\le2t_k\log(4/3)=O(k).}
\tag{CGP24}
\] This error is much smaller than \(kq\), regardless of how small the true eigenvalues are; their actual lower bound was explicitly paid for in the sample count.

The number of samples is finite with a proved bound \[
T_N=O_A(q\sqrt k\log(q+2)).
\tag{CGP25}
\] Indeed \(\log(s/r_c)\ge(s-r_c)/s=\eta/s\ge\eta\), \(D=O(q)\), CGP14 has \(\log B_N^2=o(q)\), and CGP16 gives \(\log(1/\lambda_N)=O_A(q\log(q+2))\). The values in CGP17 and the Gram CGP6 are exact real or complex quantities in this theorem. Floating-point computation would require additional certified evaluation errors; none is silently declared exact.

\subsection{5. The discarded pole rows have a proved complementary determinant cost}\label{the-discarded-pole-rows-have-a-proved-complementary-determinant-cost}

Here the word complementary refers to an exact Schur complement, not removal of those rows from the original source. The full actual row space still has dimension \(r\). Choose its coefficient- orthonormal basis with the \(t_k\) rows above first, and the remaining \(h_k\) rows second. In that basis set \[
U_N=\max\{1,V_NM_D^{2D}e^{2J_D}\},\qquad
\widetilde\lambda_N=\min\{1,(M_D^2B_E)^{-1}\}.
\tag{CGP26}
\] The full unprojected covariance lies between \(\widetilde\lambda_N I\) and \(U_NI\), by CGS14--15. The full projected covariance lies between \(\lambda_N I\) and \(U_NI\), by CGS6 and CGR11. These bounds pass to the complementary Schur minima: a lower scalar bound holds before minimization, and the lift with first coordinates zero proves the upper scalar bound. Thus, for the two rank-\(h_k\) Schur complements, every ratio of determinants has logarithm in \[
\left[h_k\log(\lambda_N/U_N),
h_k\log(U_N/\widetilde\lambda_N)\right].
\] The complete covariance determinant factors as first restriction times complementary Schur determinant. This proves \[
\boxed{|\log\det S_N-\log\det S_N^0|
\le P_N:=h_k\max\left\{
\log(U_N/\lambda_N),\log(U_N/\widetilde\lambda_N)\right\}
=O_A(q\sqrt k\log^2(q+2))=o(kq).}
\tag{CGP27}
\] Here \(S_N^0\) is formed from the two restricted covariances as displayed above. It is not asserted to equal a principal submatrix of the full \(S_N\). The exact restriction, Schur decomposition, and full scalar bounds prove CGP27 without such an identification. All quantities in \(P_N\) are specified by actual data, and \(h_k\) is the rank of the actual value/jet constraint matrix in CGR23, not an assumed generic rank.

Combining CGP24 and CGP27 at the original four cutoffs gives the finite formula for the remaining four-angle return: \[
\boxed{\left|
\log\frac{\det S_{q-1}\det S_q}{\det S_{2q-1}\det S_{2q}}
-\bigl(\mathcal A_{q-1}^{[T]}+\mathcal A_q^{[T]}
-\mathcal A_{2q-1}^{[T]}-\mathcal A_{2q}^{[T]}\bigr)
\right|
\le\sum_NP_N+8t_k\log(4/3)=o_A(kq).}
\tag{CGP28}
\] This identifies the complete finite matrix action in which any \(kq\)-scale term of the original angle return occurs. It provides an error smaller than that scale while retaining the actual analytic function \(F_z/E_A\), every lower-root projection constraint and every source degree. It does not replace that finite determinant by one, or by its rank, or by the unprojected norm. Determining its asymptotic coefficient or sign remains the mathematical calculation.

\subsection{Exact sources and reading coverage}\label{exact-sources-and-reading-coverage}

The complete KF1--55, CGR1--29 and CGS1--22 arguments were read in \nolinkurl{ACTUAL_KERNEL_FILTRATION_AND_GRAPH.md}, \nolinkurl{INVARIANT_COVARIANCE_POLE_FILTRATION.md} and \nolinkurl{INVARIANT_STRIP_ANGLE_REDUCTION.md}. The current calculation uses the exact original row and source definitions, CGR19/23--27's analytic bounds, and CGS6/14--17's full positive metric bounds. The complete sources must accompany public use until their pinned proof-file links are verified. Their equation prefixes alone do not identify a source.

The original conductor strip identities OCF4--8 are from the \href{https://github.com/KokunoYumeto/zeta-function-research-reader/blob/5992ad50c0f2a06921f1fcaded7b4e38097c0739/workbenches/splitzero-tandem/continuations/20260920-original-kernel-mixed-determinants/providers/ORIGINAL_CONDUCTOR_STRIP_FRAME.tex\#L27}{complete pinned OCF proof}, read completely and byte-verified in the parent source ledger. The monic Gamma recurrence and generating function used in CGP3, CGP6 and CGP7--8 are the Meixner--Pollaczek identities in NIST DLMF \href{https://dlmf.nist.gov/18.22.E8}{18.22.8} and \href{https://dlmf.nist.gov/18.23.E7}{18.23.7}, documented by T. H. Koornwinder, R. Wong, R. Koekoek, R. F. Swarttouw and W. P. Reinhardt. Their original TeX encodings were read and retained in \nolinkurl{graph_primary_sources/}. The classical Christoffel--Darboux identity is derived here from that exact recurrence rather than invoked as an additional estimate. The discrete Cauchy identity, its alias error and every determinant comparison used here are proved in the text above.

The independent exact checker \nolinkurl{check_original_projection_cauchy.py} passed 108 tests. They cover the Christoffel--Darboux sum and its diagonal limit, all original lower-root projection equations on an auxiliary quartet, agreement with the complete polynomial-ideal projection, discrete Cauchy aliases, and the positive relative covariance/determinant bounds. Its auxiliary rational-function fixture required 36 samples under the explicit finite sample rule. These tests do not supply an asymptotic value for the original period-dependent determinant.
